\phantomsection\addcontentsline{toc}{section}{Electric Vehicles}
\ECchapter{06}{Electric Vehicles}{Are electric cars really sustainable?}{ECBlue}

\ECObjectives{
\begin{itemize}
\item Describe the main components of a battery-electric vehicle powertrain.
\item Explain charging, regenerative braking and battery thermal management.
\item Assess an electric-vehicle project using technical and environmental criteria.
\end{itemize}}

\begin{multicols}{2}
\begin{engineeringnews}
Electric mobility is moving beyond the vehicle itself. Engineers now design complete systems
that connect cars, charging stations, buildings and electricity networks. Smart charging can
shift demand away from peak hours, while vehicle-to-grid operation may allow parked vehicles
to provide flexibility to the power system.
\end{engineeringnews}

\begin{ecarticle}
A battery-electric vehicle stores electrical energy in a high-voltage battery pack. During
acceleration, the battery supplies direct current to a power-electronics inverter. The inverter
controls the alternating current delivered to the traction motor, which converts electrical
energy into mechanical torque at the wheels. Unlike a combustion engine, an electric motor can
produce high torque from very low speed and commonly operates with high efficiency over a broad
range of conditions.

The battery is not simply a collection of cells. It also includes electrical connections,
contactors, sensors, cooling channels, structural protection and a \textbf{Battery Management
System} (BMS). The BMS estimates state of charge and state of health, balances the cells and
protects the pack against excessive current, voltage and temperature. Thermal management is
critical because low temperatures reduce available power, while high temperatures accelerate
ageing and may increase safety risks.

When the driver brakes, the inverter can operate the traction motor as a generator. Part of the
vehicle's kinetic energy is then converted back into electrical energy and stored in the battery.
This \textbf{regenerative braking} reduces energy consumption and mechanical-brake wear, although
friction brakes remain necessary for emergency stops, low-speed operation and situations in which
the battery cannot accept additional power.

Charging power depends on both the charging station and the vehicle. An on-board charger converts
alternating current from a domestic or public AC supply. Fast DC charging bypasses the on-board
charger and feeds controlled direct current to the battery. High charging power shortens stops,
but it also increases heat generation, infrastructure cost and stress on the local electricity
network.

Electric vehicles have no exhaust emissions during use, but they are not impact-free. Engineers
must consider battery manufacture, electricity generation, vehicle mass, tyre particles, lifetime
and recycling. Their environmental performance therefore depends on the complete life cycle and on
how the electricity used for charging is produced.
\end{ecarticle}

\begin{tcolorbox}[title=\sffamily\bfseries Did You Know?,colback=yellow!10,colframe=ECOrange]
An electric motor can reverse its energy flow in a fraction of a second: it drives the wheels during
acceleration and becomes a generator during regenerative braking.
\end{tcolorbox}

\columnbreak

\begin{engineeringfocus}
\textbf{Energy flow in a battery-electric vehicle}

\begin{center}
\resizebox{\linewidth}{!}{%
\begin{tikzpicture}[node distance=4mm and 5mm,font=\scriptsize,>=Latex]
\tikzset{block/.style={draw=ECBlue,fill=ECLightBlue,rounded corners,align=center,
minimum height=9mm,minimum width=18mm}}
\node[block] (bat) {battery\\pack};
\node[block,right=of bat] (inv) {inverter};
\node[block,right=of inv] (mot) {electric\\motor};
\node[block,right=of mot] (gear) {reduction\\gear};
\node[draw=ECGreen,fill=ECGreen!10,circle,right=of gear,minimum size=10mm] (wheel) {wheel};
\node[block,below=7mm of inv] (bms) {BMS};
\node[block,below=7mm of mot] (cool) {thermal\\management};
\draw[->,very thick,ECBlue] (bat) -- node[above]{DC} (inv);
\draw[->,very thick,ECBlue] (inv) -- node[above]{controlled AC} (mot);
\draw[->,very thick,ECBlue] (mot) -- (gear);
\draw[->,very thick,ECBlue] (gear) -- (wheel);
\draw[<->,ECOrange,thick] (bms) -- (bat);
\draw[<->,ECGreen,thick] (cool) -- (mot);
\draw[<->,ECGreen,thick] (cool.west) -| (bat.south);
\draw[<-,very thick,ECPurple,bend left=38] (bat.north) to node[above]{regeneration} (mot.north);
\end{tikzpicture}%
}
\end{center}

The high-voltage traction circuit is complemented by a DC/DC converter that supplies the
vehicle's low-voltage network. Isolation monitoring and contactors disconnect the battery if a
serious fault is detected.
\end{engineeringfocus}

\begin{keyfigures}
\begin{tabularx}{\linewidth}{@{}Xr@{}}
Electric motor peak efficiency & 90--97\%\\
Typical battery voltage & 350--800 V\\
Home AC charging & 2--11 kW\\
Rapid DC charging & 50--350 kW\\
Typical passenger-car battery & 40--100 kWh
\end{tabularx}
\end{keyfigures}

\begin{tcolorbox}[title=\sffamily\bfseries From supplied energy to the wheels,colback=white,colframe=ECBlue]
\begin{center}
\begin{tikzpicture}
\begin{axis}[
width=\linewidth,height=47mm,
ybar,bar width=13pt,
ymin=0,ymax=100,
ylabel={Energy reaching the wheels (kWh)},
symbolic x coords={Battery EV,Hydrogen EV,Petrol car},
xtick=data,
xticklabel style={font=\scriptsize,align=center},
yticklabel style={font=\scriptsize},
grid=major,
nodes near coords,
every node near coord/.append style={font=\scriptsize}]
\addplot[fill=ECBlue] table[x=vehicle,y=wheel_energy,col sep=comma]{data/EC06/energy_pathways.csv};
\end{axis}
\end{tikzpicture}
\end{center}
\footnotesize Illustrative conversion chain starting from 100~kWh of supplied energy. Actual
values vary with technology, driving conditions and system boundaries.
\end{tcolorbox}

\begin{vocabulary}
\begin{tabularx}{\linewidth}{@{}lX@{}}
battery pack & pack batterie\\
state of charge & état de charge\\
state of health & état de santé\\
traction motor & moteur de traction\\
regenerative braking & freinage régénératif\\
on-board charger & chargeur embarqué\\
charging point & point de recharge\\
power electronics & électronique de puissance
\end{tabularx}
\end{vocabulary}

\begin{questions}
\item What are the functions of the inverter and the traction motor?
\item Why does a battery pack require thermal management?
\item Explain how regenerative braking changes the direction of energy flow.
\item Compare AC charging with rapid DC charging.
\item Why can charging power be limited by the vehicle rather than the station?
\item Which life-cycle impacts must be considered beyond exhaust emissions?
\item Why is smart charging useful for the electricity network?
\end{questions}

\begin{engineeringchallenge}
A city operates 24 buses, each travelling 220~km per day. Compare overnight depot charging with
high-power opportunity charging at selected terminal stops. Propose a battery capacity, charging
power and operating schedule. Justify your choices using range, charging time, battery ageing,
network connection, redundancy, safety and maintenance.
\end{engineeringchallenge}

\begin{applicationexercise}
An electric bus has a usable battery capacity of 320~kWh and consumes on average
1.25~kWh/km. Estimate its theoretical range. Then calculate the energy recovered during
a 12~km downhill section if the mechanical energy available at the wheels is 18~kWh and
the regenerative chain efficiency is 72\%. Explain why the practical operating range
must remain lower than the theoretical value.
\end{applicationexercise}

\begin{speaking}
Prepare a four-minute technical briefing answering the question: ``Should all new urban cars be
electric by 2035?'' Present one technical advantage, one limitation and one policy measure that
could improve the overall transport system.
\end{speaking}
\end{multicols}

\subsection*{Sources}
\ECsource{IEA -- Global EV Outlook}{https://www.iea.org/reports/global-ev-outlook-2025}
\ECsource{U.S. Department of Energy -- Electric vehicles}{https://www.energy.gov/eere/electricvehicles/electric-vehicles}
\ECsource{Alternative Fuels Data Center -- How electric vehicles work}{https://afdc.energy.gov/vehicles/how-do-all-electric-cars-work}
